\chapter*{Introduction}
\addcontentsline{toc}{chapter}{Introduction}
\markboth{Introduction}{}

La formation pratique constitue un pilier fondamental dans le parcours acad\'emique
des \'etudiants, permettant l'acquisition de comp\'etences professionnelles et
l'application concr\`ete des connaissances th\'eoriques acquises en cours. Les stages
repr\'esentent ainsi un pont essentiel entre le monde acad\'emique et le monde
professionnel.

Au B\'enin, de plus en plus d'entreprises, \'ecoles, administrations et PME se dotent
d'un site web pour exister en ligne -- vitrine, services, communication. Mais la
s\'ecurit\'e informatique reste un domaine technique r\'eserv\'e \`a des sp\'ecialistes, et
la majorit\'e de ces structures n'ont ni les comp\'etences ni les moyens de v\'erifier
si leur site respecte les standards de s\'ecurit\'e de base. Un site web mal s\'ecuris\'e
expose son propri\'etaire, et ses visiteurs, \`a des risques concrets~: usurpation via
un site sans HTTPS, interception de donn\'ees sur une connexion non chiffr\'ee,
attaques facilit\'ees par l'absence d'en-t\^etes de s\'ecurit\'e, ou fuite d'informations
sur le serveur utilis\'e. Le probl\`eme n'est pas seulement technique~: c'est un
probl\`eme d'acc\`es \`a l'information. Le propri\'etaire d'un site ne sait g\'en\'eralement
pas si son site est vuln\'erable, ni quoi v\'erifier, ni comment corriger. Les outils
existants (Mozilla Observatory, SSL Labs) sont souvent en anglais, dispers\'es, et
pas pens\'es pour un utilisateur non-technicien qui veut simplement comprendre
\og~est-ce que mon site est s\^ur, et que dois-je faire~?~\fg{}

C'est dans ce contexte que nous avons con\c{c}u \textbf{CyberGuard B\'enin}, une
plateforme qui automatise ce diagnostic. L'utilisateur cr\'ee un compte, saisit
l'URL de son site, et la plateforme lance une s\'erie de v\'erifications non
intrusives (HTTPS, certificat SSL, en-t\^etes de s\'ecurit\'e, exposition
d'informations serveur, cookies). Le syst\`eme traduit ces r\'esultats techniques en
un score sur 100 et un niveau de risque (Bon / Moyen / Faible), accompagn\'e de
recommandations concr\`etes et compr\'ehensibles.

Ainsi, dans le cadre de notre formation, nous avons effectu\'e un stage d'un
mois au sein de la \entreprise{}, du \datedebut{} au \datefin{}, qui nous a
permis de r\'ealiser ce projet. Ce stage s'inscrit dans notre cursus acad\'emique
et constitue une \'etape cruciale de notre parcours d'apprenants.

Affect\'es \`a la Cellule de la Statistique, de la Documentation et du
Pr\'earchivage de la DGRSI, nous avons men\'e en parall\`ele de nos activit\'es
courantes le pr\'esent projet de programmation, CyberGuard B\'enin, dans le cadre
de notre m\'emoire de fin de formation en Licence Professionnelle.

Le pr\'esent rapport, structur\'e en trois parties principales, a pour objectif de
pr\'esenter le travail effectu\'e durant notre stage. Apr\`es avoir pr\'esent\'e le
d\'eroulement du stage, incluant une description de l'entreprise d'accueil, de son
environnement technologique et des missions qui nous ont \'et\'e confi\'ees, nous
abordons la description et l'analyse conceptuelle qui ont permis l'aboutissement
du projet~: l'\'etude des besoins, l'analyse fonctionnelle, ainsi que la
mod\'elisation du syst\`eme \`a travers les divers diagrammes UML (cas d'utilisation,
classes, objets, s\'equence). Enfin, nous pr\'esentons les \'etapes de
l'impl\'ementation proprement dite, les choix technologiques effectu\'es, les
difficult\'es rencontr\'ees et les solutions apport\'ees, ainsi que les r\'esultats
obtenus et les perspectives d'am\'elioration.
